% ============================================================
%  preamble.tex
% ============================================================


% ------------------------------------------------------------
%  1. Codificación, idioma y tipografía
% ------------------------------------------------------------
\usepackage[utf8]{inputenc}
\usepackage[T1]{fontenc}
\usepackage[spanish,es-tabla]{babel}

\usepackage{amsmath,amssymb}
\usepackage[sc]{mathpazo}          % Palatino (texto y matemáticas)
\linespread{1.05}                  % Recomendado para Palatino
\usepackage{microtype}
\usepackage{mathtools}
\usepackage{bm}

% Nombres en español: «Tema» en lugar de «Capítulo», «Índice» en lugar de «Índice general»
\addto\captionsspanish{%
  \renewcommand{\chaptername}{Tema}%
  \renewcommand{\contentsname}{Índice}%
}


% ------------------------------------------------------------
%  2. Geometría de página (preparada para imprimir a doble cara)
% ------------------------------------------------------------
\usepackage[a4paper,twoside,
            inner=3cm,outer=2.5cm,   % el margen interior es mayor (encuadernación)
            top=3cm,bottom=3cm,
            headheight=15pt,headsep=0.8cm]{geometry}


% ------------------------------------------------------------
%  3. Encabezados y pies de página
%     - Páginas impares (derecha):  sección ........ nº de página
%     - Páginas pares (izquierda):  nº de página ........ tema
%     - Línea horizontal bajo el encabezado, pie de página vacío
% ------------------------------------------------------------
\usepackage{fancyhdr}
\usepackage{emptypage}             % Sin encabezado en las páginas en blanco

\pagestyle{fancy}
\fancyhf{}                         % Limpia encabezado y pie
\fancyhead[LE,RO]{\small\thepage}
\fancyhead[RE]{\small\itshape\nouppercase{\leftmark}}
\fancyhead[LO]{\small\itshape\nouppercase{\rightmark}}
\renewcommand{\headrulewidth}{0.4pt}
\renewcommand{\footrulewidth}{0pt}

% Primera página de cada tema (y del índice): solo el número de página, sin línea
\fancypagestyle{plain}{%
  \fancyhf{}%
  \fancyhead[LE,RO]{\small\thepage}%
  \renewcommand{\headrulewidth}{0pt}%
}

% Qué se escribe en el encabezado
\renewcommand{\chaptermark}[1]{%
  \markboth{\chaptername\ \thechapter.\ #1}{\chaptername\ \thechapter.\ #1}}
\renewcommand{\sectionmark}[1]{\markright{\thesection.\ #1}}


% ------------------------------------------------------------
%  4. Títulos de tema y de sección
% ------------------------------------------------------------
\usepackage{titlesec}

\titleformat{\chapter}[display]
  {\normalfont\raggedright}
  {\large\scshape\chaptername\ \thechapter}
  {1ex}
  {\huge\bfseries}
  [\vspace{1.5ex}\titlerule]
\titlespacing*{\chapter}{0pt}{0pt}{4ex}

\titleformat{\section}
  {\normalfont\Large\bfseries}{\thesection}{0.8em}{}
\titlespacing*{\section}{0pt}{4ex plus 1ex minus .3ex}{2ex plus .2ex}

\titleformat{\subsection}
  {\normalfont\large\bfseries}{\thesubsection}{0.8em}{}
\titlespacing*{\subsection}{0pt}{3ex plus 1ex minus .3ex}{1.5ex plus .2ex}

% Índice: temas y secciones
\setcounter{tocdepth}{1}
\setcounter{secnumdepth}{2}


% ------------------------------------------------------------
%  5. Entornos base: Enunciado, Resolución, Nota
%     (sin colores ni recuadros; se distinguen por la tipografía)
%     Uso:  \begin{enunciado} ... \end{enunciado}
%           \begin{resolucion} ... \end{resolucion}   (sin tilde en el código)
%           \begin{nota}       ... \end{nota}
% ------------------------------------------------------------
\usepackage{amsthm}

\newtheoremstyle{enunciado}
  {1.5ex plus .5ex minus .3ex}{1.5ex plus .5ex minus .3ex}
  {\normalfont}{0pt}{\bfseries}{.}{0.6em}{}

\newtheoremstyle{resolucion}
  {1.5ex plus .5ex minus .3ex}{1.5ex plus .5ex minus .3ex}
  {\normalfont}{0pt}{\itshape}{.}{0.6em}{}

\newtheoremstyle{nota}
  {1.5ex plus .5ex minus .3ex}{1.5ex plus .5ex minus .3ex}
  {\normalfont\small}{0pt}{\scshape}{.}{0.6em}{}

\theoremstyle{enunciado}\newtheorem*{enunciado}{Enunciado}
\theoremstyle{resolucion}\newtheorem*{resolucion}{Resolución}
\theoremstyle{nota}\newtheorem*{nota}{Nota}


% ------------------------------------------------------------
%  6. Listas, figuras y enlaces
% ------------------------------------------------------------
\usepackage{enumitem}
\setlist{itemsep=0.3ex,topsep=0.8ex,parsep=0pt}
\setlist[enumerate,1]{label=(\alph*),leftmargin=2em}

\usepackage{graphicx}
\graphicspath{{imagenes/}}         % Cada tema ajusta esta ruta (ver temas/tema1.tex)
\usepackage[font=small,labelfont=bf,labelsep=period]{caption}

\usepackage[hidelinks]{hyperref}   % Enlaces internos sin colores ni recuadros


% ------------------------------------------------------------
%  7. Macros de notación (Mecánica Cuántica)
% ------------------------------------------------------------
\newcommand{\ii}{\mathrm{i}}                 % unidad imaginaria
\newcommand{\ee}{\mathrm{e}}                 % número e
\newcommand{\dd}{\mathrm{d}}                 % diferencial
\newcommand{\vect}[1]{\bm{#1}}               % vectores en negrita
\newcommand{\dv}[2]{\frac{\dd #1}{\dd #2}}   % derivada
\newcommand{\pdv}[2]{\frac{\partial #1}{\partial #2}}

\DeclareMathOperator{\Tr}{Tr}
\DeclareMathOperator{\sgn}{sgn}

% Delimitadores: \ket{\psi}, \bra{\psi}, y versiones con tamaño automático \ket*{\psi}
\DeclarePairedDelimiter{\abs}{\lvert}{\rvert}
\DeclarePairedDelimiter{\norm}{\lVert}{\rVert}
\DeclarePairedDelimiter{\bra}{\langle}{\rvert}
\DeclarePairedDelimiter{\ket}{\lvert}{\rangle}
\DeclarePairedDelimiterX{\braket}[2]{\langle}{\rangle}{#1\,\delimsize\vert\,#2}
\DeclarePairedDelimiterX{\mel}[3]{\langle}{\rangle}{#1\,\delimsize\vert\,#2\,\delimsize\vert\,#3}
\DeclarePairedDelimiterX{\ev}[1]{\langle}{\rangle}{#1}     % valor esperado
\DeclarePairedDelimiterX{\comm}[2]{[}{]}{#1,#2}            % conmutador
\DeclarePairedDelimiterX{\acomm}[2]{\{}{\}}{#1,#2}         % anticonmutador


% ------------------------------------------------------------
%  8. Portada (usa \title, \author y \date definidos en main.tex)
% ------------------------------------------------------------
\makeatletter
\renewcommand{\maketitle}{%
  \begin{titlepage}
    \setcounter{page}{1}%   La portada cuenta como página 1 (derecha)
    \centering
    \vspace*{2cm}
    {\Huge\bfseries\@title\par}
    \vspace{1.5cm}
    \includegraphics[width=0.75\textwidth,height=0.4\textheight,keepaspectratio]%
      {imagenes/portada.png}\par
    \vfill
    {\Large\@author\par}
    \vspace{0.5cm}
    {\large\@date\par}
    \vspace*{1cm}
  \end{titlepage}%
}
\makeatother